% =====================================================================
%  Thème 5 — Angle et identités remarquables — MOYEN — Exercices
% =====================================================================
\documentclass[11pt,a4paper]{article}
\usepackage[utf8]{inputenc}
\usepackage[T1]{fontenc}
\usepackage{lmodern}
\usepackage[french]{babel}
\usepackage{amsmath,amssymb,mathtools}
\usepackage{xcolor}
\usepackage{enumitem}
\usepackage{fancyhdr}
\usepackage[a4paper,margin=2.1cm,headheight=26pt,headsep=10pt]{geometry}

\definecolor{primary}{RGB}{222,70,80}
\definecolor{primarydark}{RGB}{150,30,42}

\pagestyle{fancy}\fancyhf{}
\fancyhead[L]{\small\textcolor{primarydark}{\textbf{Fiche 2} — Calcul vectoriel}}
\fancyhead[R]{\small\textcolor{primarydark}{\textsc{Thème 5 — Moyen — Exercices}}}
\fancyfoot[C]{\small\thepage}
\renewcommand{\headrulewidth}{0.8pt}

\newcommand{\vc}[1]{\overrightarrow{#1}}
\providecommand{\norm}[1]{\left\lVert #1\right\rVert}
\newcommand{\exo}[1]{\medskip\noindent{\color{primary}\bfseries\large Exercice #1.}\par\nopagebreak\smallskip}

\begin{document}
\begin{center}
{\LARGE\bfseries\color{primarydark}Angle et identités remarquables — Niveau Moyen}\\[2pt]
{\color{primary}\rule{0.55\linewidth}{1.1pt}}
\end{center}
\vspace{1mm}
{\small\itshape Objectif : mener un calcul d'angle complet, jongler entre méthode directe et identités
remarquables, et calculer un angle dans un triangle.}
\vspace{2mm}

\exo{1} \textbf{Calcul d'angle complet.}
On donne $\vec u\,(1;2)$ et $\vec v\,(4;0)$.
\begin{enumerate}[label=\textbf{\alph*)},leftmargin=2em,itemsep=2pt]
  \item Calculer $\vec u\cdot\vec v$, $\norm{\vec u}$ et $\norm{\vec v}$.
  \item En déduire $\cos\theta$ (forme exacte), puis une valeur approchée de $\theta$ au dixième
        de degré.
\end{enumerate}

\exo{2} \textbf{Angle remarquable exact.}
On donne $\vec u\,(0;2)$ et $\vec v\,(2;2\sqrt3)$. Déterminer la valeur \emph{exacte} de l'angle
$(\vec u,\vec v)$.

\exo{3} \textbf{Norme d'une somme, deux méthodes.}
On donne $\vec u\,(2;1)$ et $\vec v\,(1;3)$.
\begin{enumerate}[label=\textbf{\alph*)},leftmargin=2em,itemsep=2pt]
  \item Calculer $\norm{\vec u+\vec v}$ par la méthode directe.
  \item Retrouver $\norm{\vec u+\vec v}^2$ à l'aide de l'identité remarquable.
\end{enumerate}

\exo{4} \textbf{Retrouver un produit scalaire par une identité.}
On sait que $\norm{\vec u}=3$, $\norm{\vec v}=4$ et $\norm{\vec u+\vec v}=\sqrt{37}$.
\begin{enumerate}[label=\textbf{\alph*)},leftmargin=2em,itemsep=2pt]
  \item En utilisant $\norm{\vec u+\vec v}^2=\norm{\vec u}^2+\norm{\vec v}^2+2\,\vec u\cdot\vec v$,
        déterminer $\vec u\cdot\vec v$.
  \item En déduire $\cos(\vec u,\vec v)$, puis la mesure de l'angle.
\end{enumerate}

\exo{5} \textbf{Angle dans un triangle.}
Dans un repère orthonormé, $A\,(1;1)$, $B\,(4;2)$ et $C\,(2;5)$.
\begin{enumerate}[label=\textbf{\alph*)},leftmargin=2em,itemsep=2pt]
  \item Calculer $\vc{AB}\cdot\vc{AC}$, $\norm{\vc{AB}}$ et $\norm{\vc{AC}}$.
  \item En déduire une valeur approchée de l'angle $\widehat{BAC}$ (au dixième de degré).
\end{enumerate}

\end{document}
